% prop-logic-semantics.tex

%%%%%%%%%%%%%%%%%%%%
\begin{frame}{}
  \begin{center}
    {\Large ``再讲语义''}

    \vspace{0.30cm}
    \fig{width = 0.60\textwidth}{figs/syntax-semantics}
  \end{center}
\end{frame}
%%%%%%%%%%%%%%%%%%%%

%%%%%%%%%%%%%%%%%%%%
\begin{frame}{}
  \begin{center}
    \fbox{
      命题逻辑公式的\blue{\bf 语义}就是该公式的 \red{``真 ($T/1/\top$)''、``假 ($F/0/\bot$)''} 值。
    }

    \[
      \alpha = (B \to (A \to C)) \leftrightarrow ((B \land A) \to C)
    \]

    \pause
    \vspace{0.60cm}
    某个公式 $\alpha$ 的\red{\bf 真假值}取决于
    \vspace{0.30cm}
    \begin{columns}
      \column{0.25\textwidth}
      \column{0.50\textwidth}
        \begin{enumerate}[(1)]
          \setlength{\itemsep}{6pt}
          \item $\alpha$ 中所含命题符号的真假值; 与
          \item 逻辑联词的语义
        \end{enumerate}
      \column{0.25\textwidth}
    \end{columns}
  \end{center}
\end{frame}
%%%%%%%%%%%%%%%%%%%%

%%%%%%%%%%%%%%%%%%%%
\begin{frame}{}
  \begin{definition}[真值指派 ($v$)]
    令 $S$ 为一个命题符号的集合。\\[3pt]
    $S$ 上的一个\red{\bf 真值指派} $v$ 是一个从 $S$ 到真假值的映射
    \[
      v: S \to \set{T, F}.
    \]
  \end{definition}

  \pause
  \vspace{0.30cm}
  \[
    \alpha = (B \to (A \to C)) \leftrightarrow ((B \land A) \to C)
  \]

  \[
    v(A) = v(B) = T, \quad v(C) = F
  \]
\end{frame}
%%%%%%%%%%%%%%%%%%%%

%%%%%%%%%%%%%%%%%%%%
\begin{frame}{}
  \[
    \alpha = (B \to (A \to C)) \leftrightarrow ((B \land A) \to C)
  \]

  \[
    v(A) = v(B) = T, \quad v(C) = F
  \]

  \pause
  \[
    \red{\boxed{\overline{v}(\alpha) = ?}}
  \]
  \pause
  \[
    \overline{v}(A \to C) = ?
    \qquad \overline{v}(B \land A) = ?
    \qquad \overline{v}((B \land A) \to C) = ?
  \]

  \pause
  \begin{definition}[真值指派的扩张 ($\overline{v}$)]
    令 $S$ 为一个命题符号的集合。
    令 $\overline{S}$ 为只含有 $S$ 中命题符号的公式集。

    $S$ 上的\red{\bf 真值指派 $v$ 的扩张}是一个从 $\overline{S}$ 到真假值的映射
    \[
      \overline{v}: \overline{S} \to \set{T, F}.
    \]
  \end{definition}
\end{frame}
%%%%%%%%%%%%%%%%%%%%

%%%%%%%%%%%%%%%%%%%%
\begin{frame}{}
  \begin{center}
    使用\red{\bf 真值表}定义\red{\bf 逻辑联词的语义}

    \vspace{0.30cm}
    \begin{definition}[真值表 (Truth Table)]
      \input{tables/truth-table}
    \end{definition}

    \vspace{0.60cm}
    ``如果中国足球队夺冠, 我就好好学习''

    \pause
    \[
      (A \land B) \to B \pause \qquad \teal{v(A) = F, v(B) = T};
        \pause \qquad \brown{v(B) = F}
    \]
  \end{center}
\end{frame}
%%%%%%%%%%%%%%%%%%%%

%%%%%%%%%%%%%%%%%%%%
\begin{frame}{}
  \[
    \alpha = (B \to (A \to C)) \leftrightarrow ((B \land A) \to C)
  \]

  \[
    v(A) = v(B) = T, \quad v(C) = F
  \]

  \pause
  \[
    \overline{v}(A \to C) = F
    \qquad \overline{v}(B \to (A \to C)) = F
  \]
  \[
    \overline{v}(B \land A) = T
    \qquad \overline{v}((B \land A) \to C) = F
  \]

  \pause
  \[
    \red{\boxed{\overline{v}(\alpha) = T}}
  \]
  \[
    \alpha \text{ 在真值指派 $v$ 下的真值为}\; T
  \]

  \pause
  \begin{definition}[满足 (Satisfy)]
    如果 $\overline{v}(\alpha) = T$, 则称真值指派 $v$ \red{\bf 满足}公式 $\alpha$。
  \end{definition}
\end{frame}
%%%%%%%%%%%%%%%%%%%%

%%%%%%%%%%%%%%%%%%%%
\begin{frame}{}
  \begin{exampleblock}{课堂练习: 构造下列公式的真值表}
    \[
      (B \to (A \to C)) \leftrightarrow ((B \land A) \to C)
    \]

    \[
      (P \land Q) \land \lnot P
    \]

    \[
      \lnot P \lor Q
    \]

    \[
      (P \land Q) \lor (\lnot P \land \lnot Q)
    \]
  \end{exampleblock}
\end{frame}
%%%%%%%%%%%%%%%%%%%%

%%%%%%%%%%%%%%%%%%%%
\begin{frame}{}
  \begin{definition}[可满足的 (Satisfiable)]
    如果\blue{\bf 存在}某个真值指派满足公式 $\alpha$,
    则称 $\alpha$ 是\blue{\bf 可满足的}。
  \end{definition}

  \[
    \lnot P \lor Q \text{ 是{可满足的}}
  \]

  \[
    \alpha = (B \to (A \to C)) \leftrightarrow ((B \land A) \to C)
    \text{ 是{可满足的}}
  \]

  \[
    (P \land Q) \land \lnot P \text{ 是\red{不可满足的} (unsatisfiable)}
  \]
\end{frame}
%%%%%%%%%%%%%%%%%%%%

%%%%%%%%%%%%%%%%%%%%
\begin{frame}{}
  \begin{definition}[命题逻辑公式的可满足性问题 (Satisfiability; \red{\bf SAT})]
    任给一个命题逻辑公式 $\alpha$, $\alpha$ 是可满足的吗?
  \end{definition}

  \pause
  \vspace{0.50cm}
  \fig{width = 0.50\textwidth}{figs/millionaire}
\end{frame}
%%%%%%%%%%%%%%%%%%%%

%%%%%%%%%%%%%%%%%%%%
\begin{frame}{}
  \begin{definition}[重言式/永真式 (Tautology)]
    若公式 $\alpha$ 在\blue{\bf 所有}真值指派下均为真, 则称 $\alpha$ 为\red{\bf 重言式},
    记为 $\models \alpha$。
  \end{definition}

  \[
    (B \to (A \to C)) \leftrightarrow ((B \land A) \to C) \text{ 是\red{重言式}}
  \]

  \[
    (\alpha \to (\beta \to \gamma)) \leftrightarrow ((\alpha \land \beta) \to \gamma)
  \]

  \pause
  \[
    \red{\boxed{(\alpha \to \beta) \leftrightarrow (\lnot \alpha \lor \beta)}}
  \]
\end{frame}
%%%%%%%%%%%%%%%%%%%%

%%%%%%%%%%%%%%%%%%%%
\begin{frame}{}
  \begin{center}
    \fig{width = 0.30\textwidth}{figs/David-Hilbert}
    \teal{David Hilbert (23 January 1862 - 14 February 1943)}
  \end{center}

  \begin{enumerate}[(1)]
    \item
      \[
        \alpha \to (\beta \to \alpha)
      \]
    \item
      \[
        (\alpha \to (\beta \to \gamma)) \to ((\alpha \to \beta) \to (\alpha \to \gamma))
      \]
    \item
      \[
        (\lnot \beta \to \lnot \alpha) \to ((\lnot \beta \to \alpha) \to \beta)
      \]
  \end{enumerate}
\end{frame}
%%%%%%%%%%%%%%%%%%%%

%%%%%%%%%%%%%%%%%%%%
\begin{frame}{}
  \begin{definition}[矛盾式/永假式 (Contradiction)]
    若公式 $\alpha$ 在\blue{\bf 所有}真值指派下均为假, 则称 $\alpha$ 为\red{\bf 矛盾式}。
  \end{definition}

  \[
    (P \land Q) \land \lnot P \text{ 是\red{矛盾式 (不可满足的)}}
  \]
\end{frame}
%%%%%%%%%%%%%%%%%%%%

%%%%%%%%%%%%%%%%%%%%
\begin{frame}{}
  \[
    \boxed{\models \alpha}
  \]

  \begin{definition}[重言式/永真式 (Tautology)]
    若公式 $\alpha$ 在\blue{\bf 所有}真值指派下均为真, 则称 $\alpha$ 为\red{\bf 重言式},
    记为 $\models \alpha$。
  \end{definition}

  \pause
  \[
    \blue{\set{\beta, \gamma, \tau}} \models \purple{\alpha}
  \]
  \begin{center}
    仅要求 \purple{\bf 结论} $\alpha$ 在 $\beta$, $\gamma$, $\tau$ 等\blue{\bf 前提条件}都为真时为真。
  \end{center}

  \pause
  \[
    \boxed{\blue{\Sigma} \models \purple{\alpha}}
  \]
  \begin{center}
    仅要求 \purple{\bf 结论} $\alpha$ 在公式集 $\Sigma$ 中所有\blue{\bf 前提条件}公式都为真时为真。
  \end{center}

  \pause
  \[
    \Sigma = \set{\alpha} \text{ 只含有一个公式时},\;
    \set{\alpha} \models \beta \text{ 简记为 } \;\alpha \models \beta
  \]

  \pause
  \[
    \Sigma = \emptyset \text{ 时},\;
    \emptyset \models \alpha \text{ 简记为 } \models \alpha
  \]
\end{frame}
%%%%%%%%%%%%%%%%%%%%

%%%%%%%%%%%%%%%%%%%%
\begin{frame}{}
  \begin{definition}[重言蕴含 (Tautologically Implies)]
    \begin{center}
      设 $\Sigma$ 为一个公式集。\\[3pt]
      如果\blue{\bf 每一个}满足 $\Sigma$ 中\blue{\bf 所有}公式的真值指派都满足 $\alpha$, \\[3pt]
      则称 $\Sigma$ \red{\bf 重言蕴含}公式 $\alpha$,
      记为 $\Sigma \models \alpha$。
    \end{center}
  \end{definition}

  \pause
  \begin{exampleblock}{小小的练习一下}
    \[
      \set{\alpha \land \beta} \models \alpha
    \]

    \pause
    \[
      \set{\alpha \to \beta, \alpha} \models \beta
    \]

    \pause
    \[
      \set{\alpha, \lnot \alpha} \models \beta
    \]

    \pause
    \[
      \alpha \models \beta \to \alpha
    \]
  \end{exampleblock}
\end{frame}
%%%%%%%%%%%%%%%%%%%%

%%%%%%%%%%%%%%%%%%%%
% \begin{frame}{}
%   \begin{theorem}
%     请证明: $\Sigma \cup \set{\alpha} \models \beta$
%     \blue{当且仅当} $\Sigma \models (\alpha \to \beta)$。
%   \end{theorem}
%
%   \pause
%   \[
%     v_{\Sigma}: \text{ 满足 $\Sigma$ 中所有公式的任一真值指派}
%   \]
%
%   \[
%     v_{\Sigma \cup \set{\alpha}}:
%     \text{ 满足 $\Sigma \cup \set{\alpha}$ 中所有公式的任一真值指派}
%   \]
%
%   \pause
%   \begin{lemma}
%     如果 $\overline{v_{\Sigma \cup \set{\alpha}}}(\beta) = T$,
%     则 $\overline{v_{\Sigma}}(\alpha \to \beta) = T$。
%   \end{lemma}
%
%   \pause
%   \vspace{0.30cm}
%   \begin{lemma}
%     如果 $\overline{v_{\Sigma}}(\alpha \to \beta) = T$,
%     则 $\overline{v_{\Sigma \cup \set{\alpha}}}(\beta) = T$。
%   \end{lemma}
% \end{frame}
%%%%%%%%%%%%%%%%%%%%

%%%%%%%%%%%%%%%%%%%%
\begin{frame}{}
  \begin{definition}[重言等价 (Tautologically Equivalent)]
    如果 $\alpha \models \beta$ 且 $\beta \models \alpha$,
    则称 $\alpha$ 与 $\beta$ \red{\bf 重言等价},
    记为 $\alpha \equiv \beta$ (或 $\alpha \Leftrightarrow \beta$)。
  \end{definition}

  \[
    (B \to (A \to C)) \leftrightarrow ((B \land A) \to C)
  \]

  \[
    B \to (A \to C) \equiv (B \land A) \to C
  \]

  \pause
  \[
    \text{注意}\; \teal{\alpha \Leftrightarrow \beta \;(\text{语义})} \text{ 与}\;
    \purple{\alpha \leftrightarrow \beta \;(\text{语法})}
    \text{ 的区别与联系}
  \]

  \pause
  \[
    \boxed{
      \teal{\alpha \Leftrightarrow \beta} \text{ 成立}
      \text{ 当且仅当}\; \purple{\alpha \leftrightarrow \beta} \text{ 是重言式}}
  \]

  \pause
  \[
    \text{如果已知}\; \purple{\alpha \leftrightarrow \beta} \text{ 是重言式,}\;
    \text{则可将}\; \alpha \text{ 替换成\teal{\bf 语义上等价的}}\; \beta\text{。}
  \]
\end{frame}
%%%%%%%%%%%%%%%%%%%%

%%%%%%%%%%%%%%%%%%%%
\begin{frame}{}
  \begin{exampleblock}{重要的重言式举例}
    \begin{description}[<+->][结合律:]
      \item[幂等律:]
        \[
          (\alpha \land \alpha) \leftrightarrow \alpha
        \]
        \[
          (\alpha \lor \alpha) \leftrightarrow \alpha
        \]
      \item[交换律:]
        \[
          (\alpha \land \beta) \leftrightarrow (\beta \land \alpha)
        \]
        \[
          (\alpha \lor \beta) \leftrightarrow (\beta \lor \alpha)
        \]
      \item[结合律:]
        \[
          ((\alpha \land \beta) \land \gamma) \leftrightarrow (\alpha \land (\beta \land \gamma))
        \]
        \[
          ((\alpha \lor \beta) \lor \gamma) \leftrightarrow (\alpha \lor (\beta \lor \gamma))
        \]
    \end{description}
  \end{exampleblock}
\end{frame}
%%%%%%%%%%%%%%%%%%%%

%%%%%%%%%%%%%%%%%%%%
\begin{frame}{}
  \begin{exampleblock}{重要的重言式举例}
    \begin{description}[<+->][分配律:]
      \item[吸收律:]
        \[
          (\alpha \land (\alpha \lor \beta)) \leftrightarrow \alpha
        \]
        \[
          (\alpha \lor (\alpha \land \beta)) \leftrightarrow \alpha
        \]
      \item[分配律:]
        \[
          (\alpha \land (\beta \lor \gamma)) \leftrightarrow ((\alpha \land \beta) \lor (\alpha \land \gamma))
        \]
        \[
          (\alpha \lor (\beta \land \gamma)) \leftrightarrow ((\alpha \lor \beta) \land (\alpha \lor \gamma))
        \]
    \end{description}
  \end{exampleblock}
\end{frame}
%%%%%%%%%%%%%%%%%%%%

%%%%%%%%%%%%%%%%%%%%
\begin{frame}{}
  \begin{exampleblock}{重要的重言式举例}
    \begin{description}[<+->][条件重言式:]
      \item[条件重言式:]
        \[
          \red{\boxed{(\alpha \to \beta) \leftrightarrow (\lnot \alpha \lor \beta)}}
        \]
    \end{description}
  \end{exampleblock}
\end{frame}
%%%%%%%%%%%%%%%%%%%%

%%%%%%%%%%%%%%%%%%%%
\begin{frame}{}
  \begin{exampleblock}{重要的重言式举例}
    \begin{description}[<+->][德摩根 (De Morgan) 律:]
      \item[德摩根 (De Morgan) 律:]
        \[
          \lnot (\alpha \land \beta) \leftrightarrow (\lnot \alpha \lor \lnot \beta)
        \]
        \[
          \lnot (\alpha \lor \beta) \leftrightarrow (\lnot \alpha \land \lnot \beta)
        \]
    \end{description}
  \end{exampleblock}
  \fig{width = 0.25\textwidth}{figs/Augustus-De-Morgan}
\end{frame}
%%%%%%%%%%%%%%%%%%%%

%%%%%%%%%%%%%%%%%%%%
\begin{frame}{}
  \begin{exampleblock}{重要的重言式举例}
    \begin{description}[<+->][双重否定律:]
      \setlength{\itemsep}{8pt}
      \item[双重否定律:]
        \[
          \lnot \lnot \alpha \leftrightarrow \alpha
        \]
      \item[排中律:]
        \[
          \alpha \lor (\lnot \alpha)
        \]
      \item[矛盾律:]
        \[
          \lnot (\alpha \land \lnot \alpha)
        \]
      \item[逆否命题:]
        \[
          (\alpha \to \beta) \leftrightarrow (\lnot \beta \to \lnot \alpha)
        \]
    \end{description}
  \end{exampleblock}
\end{frame}
%%%%%%%%%%%%%%%%%%%%

%%%%%%%%%%%%%%%%%%%%
\begin{frame}{}
  \begin{exampleblock}{请证明以下重言等价}
    \[
      (\alpha \lor \beta) \to \gamma \equiv (\alpha \to \gamma) \land (\beta \to \gamma)
    \]
  \end{exampleblock}
\end{frame}
%%%%%%%%%%%%%%%%%%%%

%%%%%%%%%%%%%%%%%%%%
\begin{frame}{}
  \begin{exampleblock}{请证明以下公式是重言式 (提示: 使用前面提到的重言式)}
    \[
      ((\alpha \to \beta) \land \alpha) \to \beta
    \]

    \pause
    \[
      ((\alpha \to \beta) \land \lnot \beta) \to \lnot \alpha
    \]

    \pause
    \[
      ((\alpha \lor \beta) \land \lnot \alpha) \to \beta
    \]
  \end{exampleblock}
\end{frame}
%%%%%%%%%%%%%%%%%%%%

%%%%%%%%%%%%%%%%%%%%
% \begin{frame}{}
%   \begin{exampleblock}{使用命题逻辑进行推理}
%     某公司要从赵、钱、孙、李、吴5名员工中选派某些人出国考察。\\
%     由于某些不可描述的原因, 选派要求如下:
%
%     \begin{columns}
%       \column{0.50\textwidth}
%         \begin{enumerate}[(1)]
%           \item 若赵去, 钱也去;
%           \item 李、吴两人中必有一人去;
%           \item 钱、孙两人中去且仅去一人;
%           \item 孙、李两人同去或同不去;
%           \item 若吴去, 则赵、钱也去;
%           \item 只有孙去, 赵才会去。
%         \end{enumerate}
%       \column{0.50\textwidth}
%         \begin{enumerate}[(1)]
%           \item $\teal{Z \to Q}$;
%           \item $\teal{L \lor W}$;
%           \item $\teal{(Q \land \lnot S) \lor (S \land \lnot Q)}$;
%           \item $\teal{(S \land L) \lor (\lnot S \land \lnot L)}$;
%           \item $\teal{W \to Z \land Q}$;
%           \item $\teal{Z \to S}$。
%         \end{enumerate}
%     \end{columns}
%
%     \vspace{0.50cm}
%     \blue{请使用形式化推理的方法帮该公司判断应选哪些人出国考察。}
%   \end{exampleblock}
% \end{frame}
%%%%%%%%%%%%%%%%%%%%

% %%%%%%%%%%%%%%%%%%%%
% \begin{frame}{}
%   \begin{definition}[合取范式 (Conjunctive Normal Form)]
%     我们称公式 $\alpha$ 是\red{\bf 合取范式}, 如果它形如
%     \[
%       \alpha = \beta_{1} \land \beta_{2} \land \dots \land \beta_{k},
%     \]
%     其中, 每个 $\beta_{i}$ 都形如
%     \[
%       \beta_{i} = \beta_{i1} \lor \beta_{i2} \lor \dots \lor \beta_{in},
%     \]
%     并且 $\beta_{ij}$ 或是一个命题符号, 或者命题符号的否定。
%   \end{definition}
%
%   \[
%     (P \lor \lnot Q \lor R) \;\red{\land}\; (\lnot P \lor Q) \;\red{\land}\; \lnot Q
%   \]
% \end{frame}
% %%%%%%%%%%%%%%%%%%%%
%
% %%%%%%%%%%%%%%%%%%%%
% \begin{frame}{}
%   \begin{exampleblock}{求下列公式的合取范式}
%     \[
%       (P \land (Q \to R)) \to S
%     \]
%
%     \[
%       \lnot (P \lor Q) \leftrightarrow (P \land Q)
%     \]
%   \end{exampleblock}
% \end{frame}
% %%%%%%%%%%%%%%%%%%%%
%
% %%%%%%%%%%%%%%%%%%%%
% \begin{frame}{}
%   \begin{definition}[析取范式 (Disjunctive Normal Form)]
%     我们称公式 $\alpha$ 是\red{\bf 析取范式}, 如果它形如
%     \[
%       \alpha = \beta_{1} \lor \beta_{2} \lor \dots \lor \beta_{k},
%     \]
%     其中, 每个 $\beta_{i}$ 都形如
%     \[
%       \beta_{i} = \beta_{i1} \land \beta_{i2} \land \dots \land \beta_{in},
%     \]
%     并且 $\beta_{ij}$ 或是一个命题符号, 或者命题符号的否定。
%   \end{definition}
% 
%   \[
%     (P \land \lnot Q \land R) \;\red{\lor}\; (\lnot P \land Q) \;\red{\lor}\; \lnot Q
%   \]
% \end{frame}
% %%%%%%%%%%%%%%%%%%%%
% 
% %%%%%%%%%%%%%%%%%%%%
% \begin{frame}{}
%   \begin{exampleblock}{求下列公式的析取范式}
%     \[
%       (P \land (Q \to R)) \to S
%     \]
%
%     \[
%       \lnot (P \lor Q) \leftrightarrow (P \land Q)
%     \]
%   \end{exampleblock}
% \end{frame}
% %%%%%%%%%%%%%%%%%%%%
%
% %%%%%%%%%%%%%%%%%%%%
% \begin{frame}{}
%   \begin{exampleblock}{求合取范式与析取范式的方法}
%     \begin{enumerate}[(1)]
%       \setlength{\itemsep}{8pt}
%       \item 先将公式中的联结符号化归成 $\lnot$, $\land$ 与 $\lor$;
%       \item 再使用 De Morgan 律将 $\lnot$ 移到各个命题变元之前 (\blue{``否定深入''});
%       \item 最后使用结合律、分配律将公式化归成合取范式或析取范式。
%     \end{enumerate}
%   \end{exampleblock}
% \end{frame}
% %%%%%%%%%%%%%%%%%%%%